% ==============================================================================
% PAGE 33: TRACKING ACTIVATION & ADVANTAGES (Numbered Page 25)
% ==============================================================================
\noindent
{\fontsize{10.2}{14.5}\selectfont
\textbf{$\blacklozenge$ Motor Activation:}
\begin{itemize}[leftmargin=6mm, itemsep=1.5mm, label=\textbullet]
  \item The comparator or controller sends a signal to the motor driver.
  \item The PMDC motor rotates the dish toward the brighter side.
\end{itemize}

\vspace{2mm}
\textbf{$\blacklozenge$ Balancing:}
\begin{itemize}[leftmargin=6mm, itemsep=1.5mm, label=\textbullet]
  \item When both sensors receive equal light again, the motor stops.
  \item The dish stays aligned with the sun automatically.
\end{itemize}

\vspace{2mm}
\textbf{$\blacklozenge$ End of Day:}
\begin{itemize}[leftmargin=6mm, itemsep=1.5mm, label=\textbullet]
  \item The dish remains in the West position at sunset.
  \item It can be manually or automatically reset to the East in the morning.
\end{itemize}
}

\vspace{4mm}
\subsubsection*{5. Functions of Tracking Mechanism}
\begin{itemize}[leftmargin=8mm, itemsep=2mm, label=\textbullet]
  \item Continuously aligns the dish with the sun's position.
  \item Ensures maximum solar radiation is focused on the receiver.
  \item Eliminates need for manual adjustment.
  \item Increases water heating and overall system efficiency.
  \item Operates using solar power, making it self-sufficient.
\end{itemize}

\vspace{4mm}
\subsubsection*{6. Advantages}
\begin{itemize}[leftmargin=8mm, itemsep=2mm, label=\textbullet]
  \item Simple and cost-effective compared to dual-axis systems.
  \item Reduces human effort.
  \item Compact and energy-efficient.
  \item Ideal for small- to medium-size parabolic dish systems.
  \item Improves output by 20--30\% over fixed-position setup.
\end{itemize}
\clearpage
